%% kapa: exam=1 round=24 year=2013 period=1 form=A subject=economics subject_ko=경제원론 seq=40 number=80
%% source: Q-Net 감정평가사 기출문제 (cst00309) artlSeq=5073802 file=extracted/1008/5073802/2013년 제24회 감정평가사 1차시험 기출문제/1교시(A형).hwp
%% conversion: kordoc 4.24.51 -> tools/gen1.py (auto); review=proofread-sonnet; flags=none
\begin{question}{80}
다음 모형에 기초하여 분석할 때, 물가수준-국민소득 평면에서 총수요곡선이 우하향하는 이유를 가장 적절하게 설명한 것은? (단, $Y$ 는 소득, $C$ 는 소비, $I$ 는 투자, $r$ 은 실질 이자율, $T$ 는 조세, $G$ 는 정부지출, $P$ 는 물가수준, $\pi^{ e }$ 는 예상물가상승률, $\beta_{ i } , i=0, \cdots ,5,$는 양의 계수, $M^{ d }$ 는 명목 화폐수요, $M^{ s }$는 명목 화폐공급이며, $T_{ 0 }$, $G_{ 0 }$, $\overline{ M }$, $\pi_{ 0 }^{ e }$은 상수이다.)
\begin{framedtext}
$C= \beta_{ 0 } + \beta_{ 1 } (Y-T)$ \par
$I= \beta_{ 2 } - \beta_{ 3 } r$ \par
$T=T_{ 0 }$, $G=G_{ 0 }$ \par
$Y=C+I+G$(생산물시장의 균형) \par
$\frac{ M^{ d } } { P } = \beta_{ 4 } Y- \beta_{ 5 } (r+ \pi^{ e } )$ \par
$M^{ s } = \overline{ M }$ \par
$M^{ d } =M^{ s }$(화폐시장의 균형) \par
$\pi^{ e } = \pi_{ 0 }^{ e }$
\end{framedtext}
\begin{choices}
  \item 물가수준이 하락하면 실질 자산의 가치가 증가하여 소비지출이 증가한다.
  \item 물가수준이 하락하면 실질 화폐공급의 증가로 이자율이 하락하여 투자가 증가한다.
  \item 물가수준이 하락하면 해외로 자본유출이 증가하여 순수출이 증가한다.
  \item 물가수준의 변화와 무관하게 구축효과로 인해 총수요가 감소한다.
  \item 물가수준의 변화와 무관하게 피구효과로 인해 총수요가 감소한다.
\end{choices}
\end{question}
\answer{2}
